% =====================================================================
%  CARNET D'ENTRAÎNEMENT — Fiche 12 (Chimie organique)
%  THÈME 2 — Reconnaissance et classification des groupements fonctionnels
%  NIVEAU FACILE — CORRIGÉ
% =====================================================================
\documentclass[11pt,a4paper]{article}
\usepackage[utf8]{inputenc}
\usepackage[T1]{fontenc}
\usepackage{lmodern}
\usepackage[french]{babel}
\usepackage{amsmath,amssymb,mathtools}
\providecommand{\degree}{^\circ}
\usepackage{icomma}
\usepackage{siunitx}
\sisetup{output-decimal-marker={,},per-mode=symbol}
\usepackage[version=4]{mhchem}
\usepackage{chemfig}
\setchemfig{atom sep=2em,bond length=0.9em,cram width=2.5pt}
\usepackage{xcolor}
\usepackage{array}
\usepackage{booktabs}
\usepackage{tabularx}
\usepackage[most]{tcolorbox}
\usepackage{enumitem}
\usepackage{tikz}
\usetikzlibrary{arrows.meta,positioning,calc}
\usepackage{fancyhdr}
\usepackage{titlesec}
\usepackage[hidelinks]{hyperref}
\usepackage[a4paper,margin=2.1cm,headheight=26pt,headsep=9pt]{geometry}

\definecolor{primary}{RGB}{150,98,162}
\definecolor{primarydark}{RGB}{92,52,110}
\definecolor{excol}{RGB}{0,135,90}
\definecolor{attncol}{RGB}{200,40,55}
\definecolor{methcol}{RGB}{90,90,100}

\newcommand{\runtitle}{Thème 2 — Facile — Corrigé}
\pagestyle{fancy}\fancyhf{}
\fancyhead[L]{\small\textcolor{primarydark}{\textbf{Fiche 12} — Chimie organique}}
\fancyhead[R]{\small\textcolor{primarydark}{\runtitle}}
\fancyfoot[C]{\small\thepage}
\renewcommand{\headrulewidth}{0.8pt}
\renewcommand{\headrule}{\hbox to\headwidth{\color{primary}\leaders\hrule height \headrulewidth\hfill}}
\setlength{\parindent}{0pt}\setlength{\parskip}{3pt}

\tcbset{boxbase/.style={enhanced,breakable,boxrule=0.9pt,arc=2.5pt,
   left=3mm,right=3mm,top=2mm,bottom=2mm,fonttitle=\bfseries,coltitle=white,
   toptitle=0.5mm,bottomtitle=0.5mm}}
\newtcolorbox[auto counter]{corr}[1]{boxbase,colback=excol!5,colframe=excol,
   title={\textsc{Corrigé}~\thetcbcounter\ \textmd{--- \itshape #1}}}
\newcommand{\rep}[1]{\textcolor{primarydark}{\textbf{#1}}}

\newcommand{\banner}[2]{%
\begin{tcolorbox}[enhanced,colback=excol,colframe=excol,arc=5pt,boxrule=0pt,
   left=5mm,right=5mm,top=2.5mm,bottom=2.5mm,coltext=white]
{\large\bfseries #1}\\[1pt]{\small #2}
\end{tcolorbox}\vspace{2pt}}

\begin{document}

\banner{Thème 2 — Reconnaissance des groupements fonctionnels}{Niveau \textbf{Facile} — Corrigé détaillé \quad$\vert$\quad Fiche 12, Chimie organique}

% ---------------------------------------------------------------------
\begin{corr}{nommer la fonction de chaque molécule}
\begin{itemize}[leftmargin=1.6em,itemsep=1pt]
  \item \textbf{(A)} $\ce{CH3-OH}$ : O--H sur un C saturé $\Rightarrow$ \rep{alcool} (méthanol) ;
  \item \textbf{(B)} $\ce{CH3-CO-CH3}$ : carbonyle entre deux C $\Rightarrow$ \rep{cétone} (propanone) ;
  \item \textbf{(C)} $\ce{CH3-CHO}$ : carbonyle en bout de chaîne (porte un H) $\Rightarrow$ \rep{aldéhyde} (éthanal) ;
  \item \textbf{(D)} $\ce{CH3-NH2}$ : azote sur C saturé, sans carbonyle $\Rightarrow$ \rep{amine} (méthylamine) ;
  \item \textbf{(E)} $\ce{CH3-COOH}$ : carbonyle portant un O--H $\Rightarrow$ \rep{acide carboxylique} (acide éthanoïque).
\end{itemize}
\end{corr}

% ---------------------------------------------------------------------
\begin{corr}{aldéhyde ou cétone ?}
On regarde ce qui entoure le carbonyle $\ce{C=O}$ :
\begin{itemize}[leftmargin=1.6em,itemsep=1pt]
  \item \textbf{(A)} le $\ce{C=O}$ est en \emph{bout} de chaîne : il porte un \textbf{H} d'un côté, un C de l'autre $\Rightarrow$ \rep{aldéhyde} (propanal) ;
  \item \textbf{(B)} le $\ce{C=O}$ est encadré par \textbf{deux carbones} $\Rightarrow$ \rep{cétone} (propanone).
\end{itemize}
Même groupe carbonyle, deux fonctions différentes : \emph{tout est dans le voisinage}.
\end{corr}

% ---------------------------------------------------------------------
\begin{corr}{alcool ou éther ?}
\begin{itemize}[leftmargin=1.6em,itemsep=1pt]
  \item \textbf{(A)} $\ce{CH3CH2-OH}$ : l'oxygène porte un H (O--H terminal) $\Rightarrow$ \rep{alcool} (éthanol) ;
  \item \textbf{(B)} $\ce{CH3-O-CH3}$ : l'oxygène ponte deux carbones $\Rightarrow$ \rep{éther} (diméthyléther).
\end{itemize}
Les deux ont la formule brute $\ce{C2H6O}$ mais des \emph{fonctions} différentes : ce sont des \rep{isomères de fonction} (mêmes atomes, propriétés chimiques très différentes).
\end{corr}

% ---------------------------------------------------------------------
\begin{corr}{amine ou amide ?}
\begin{itemize}[leftmargin=1.6em,itemsep=1pt]
  \item \textbf{(A)} $\ce{CH3-NH2}$ : l'azote est lié à un carbone \emph{saturé}, \textbf{sans} carbonyle voisin $\Rightarrow$ \rep{amine} (primaire) ;
  \item \textbf{(B)} $\ce{CH3-CO-NH2}$ : l'azote est lié à un \textbf{carbonyle} $\ce{C=O}$ $\Rightarrow$ \rep{amide} (éthanamide).
\end{itemize}
La présence (ou non) d'un $\ce{C=O}$ accolé à l'azote fait toute la différence.
\end{corr}

% ---------------------------------------------------------------------
\begin{corr}{classer un alcool (I / II / III)}
Dans le butan-2-ol, le carbone porteur du $\ce{-OH}$ est lié à \textbf{deux} autres carbones (le $\ce{CH3}$ à sa gauche et le $\ce{CH2}$ à sa droite). Un C--OH lié à 2 autres C $\Rightarrow$ \rep{alcool secondaire}.
\end{corr}

\end{document}
